\documentclass[a4paper]{book}

%------------------------------------------------
% Idioma y codificación
%------------------------------------------------

\usepackage[english]{babel}
\usepackage[T1]{fontenc}
\usepackage[utf8]{inputenc}

%------------------------------------------------
% Matemáticas
%------------------------------------------------

\usepackage{amsmath,amssymb,amsthm,mathtools}
\usepackage{mathrsfs}

%------------------------------------------------
% Tipografía
%------------------------------------------------

\usepackage{microtype}
\usepackage{lmodern}

%------------------------------------------------
% Márgenes
%------------------------------------------------

\usepackage[
    margin=3cm
]{geometry}

%------------------------------------------------
% Hipervínculos
%------------------------------------------------

\usepackage{hyperref}

%------------------------------------------------
% Estilo de página
%------------------------------------------------

\pagestyle{plain}

%------------------------------------------------
% Espaciado
%------------------------------------------------

\setlength{\parindent}{0pt}
\setlength{\parskip}{6pt}

%------------------------------------------------
% Teoremas
%------------------------------------------------

\numberwithin{equation}{section}

\theoremstyle{definition}
\newtheorem{theorem}{Theorem}[section]
\newtheorem{lemma}{Lemma}[section]
\newtheorem{proposition}{Proposition}[section]
\newtheorem{corollary}{Corollary}[section]
\newtheorem{definition}{Definition}[section]
\newtheorem{example}{Example}[section]
\newtheorem*{remark}{Remark}

%------------------------------------------------
% Operadores
%------------------------------------------------

\DeclareMathOperator{\im}{im}
\DeclareMathOperator{\interior}{int}
\DeclareMathOperator{\ext}{ext}
\DeclareMathOperator{\cl}{cl}
\DeclareMathOperator{\bd}{bd}
\DeclareMathOperator{\fr}{fr}

%------------------------------------------------
% Comandos
%------------------------------------------------

\newcommand{\pr}[1]{\left(#1\right)}
\newcommand{\br}[1]{\left[#1\right]}
\newcommand{\pbr}[1]{\left(#1\right]} % chktex 9
\newcommand{\bpr}[1]{\left[#1\right)} % chktex 9
\newcommand{\set}[1]{\left\{#1\right\}}
\newcommand{\bars}[1]{\left|#1\right|}
\newcommand{\sys}[1]{\left\{#1\right.}
\newcommand{\rest}[2]{\left.#1\right|_{#2}}
\newcommand{\Leb}[3]{\mathcal{L}^1_{#3}\pr{#1,#2}}
\newcommand{\LebR}[3]{\mathcal{L}^1_{#3}\pr{#1,#2;\mathbb{R}}}
\newcommand{\LebC}[3]{\mathcal{L}^1_{#3}\pr{#1,#2;\mathbb{C}}}


\begin{document}
  \tableofcontents

  \part{General Topology}
  \chapter{Introduction to General Topology}
\chapter{Limits}

\section{Sequences}

\begin{definition}
  Let $X$ be a set. A \textit{sequence in $X$} is a map $\varphi:\mathbb{N}\rightarrow X$. The map $\varphi$ will be denoted by $\set{x_n}_{n\in\mathbb{N}}$ where $x_n:=\varphi\pr{x_n}$ for all $n\in\mathbb{N}$. In order to emphasize that $\set{x_n}_{n\in\mathbb{N}}$ is a sequence in $X$, we will write $\set{x_n}_{n\in\mathbb{N}}\subseteq X$.
\end{definition}

\begin{definition}
  Let $\pr{X,\mathcal{T}}$ be a topological space. A sequence \emph{$\set{x_n}_{n\in\mathbb{N}}\subseteq X$ converges to a point $x\in X$ (with respect to $\pr{X,\mathcal{T}}$)} if for every neighborhood $N\in\mathcal{N}_x\pr{X,\mathcal{T}}$, there exists $n_0=n_0\pr{N,x}\in\mathbb{N}$ such that $x_n\in N$ for all $n\in\mathbb{N}$ with $n\geq n_0$. 
\end{definition}

\begin{remark}
  Let $\pr{X,\mathcal{T}}$ be a topological space. If \emph{$\set{x_n}_{n\in\mathbb{N}}\subseteq X$} is a sequence converging to $x\in X$, we say that \textit{$x$ is a limit of $\set{x_n}_{n\in\mathbb{N}}$}. Equivalently, we say that \emph{the limit of $\set{x_n}_{n\in\mathbb{N}}$ as $n\to\infty$ is $x$}. This is denoted it by $\lim_{n\to\infty}{x_n}=x$ or $x_n\stackrel{n\to\infty}{\rightarrow}x$. 
  
  In general, a sequence in a topological space may have more than one limit.
\end{remark}


\begin{proposition}
  Let $\pr{X,\mathcal{T}}$ be a topological space with basis $\mathcal{B}\subseteq\mathcal{T}$ for $\mathcal{T}$, $\set{x_n}_{n\in\mathbb{N}}\subseteq X$ be a sequence and $x\in X$. The following are equivalent:

  \begin{enumerate}
    \item $\set{x_n}_{n\in\mathbb{N}}$ converges to $x$.
    \item For every open set $G\in\mathcal{T}$, there exists $n_0=n_0\pr{G,x}\in\mathbb{N}$ such that $x_n\in G$ for all $n\in\mathbb{N}$ with $n\geq n_0$.
    \item For every basis set $B\in\mathcal{T}$, there exists $n_0=n_0\pr{B,x}\in\mathbb{N}$ such that $x_n\in B$ for all $n\in\mathbb{N}$ with $n\geq n_0$.
  \end{enumerate}
  \begin{proof}[\rm\textbf{Proof}]
    \hfill
    \begin{itemize}
      \item[\boxed{1\Rightarrow2}] Trivial.
      \item[\boxed{2\Rightarrow3}] Trivial.
      \item[\boxed{3\Rightarrow1}] Let $N\in\mathcal{N}_x\pr{X,\mathcal{T}}$. Then there exists $B=B\pr{x,N}\in\mathcal{B}$ such that $x\in B\subseteq N$. Thus there exists $n_0=n_0\pr{B,x}\in\mathbb{N}$ such that $x_n\in B\subseteq N$ for all $n\in\mathbb{N}$ with $n\geq n_0$. Therefore, $\set{x_n}_{n\in\mathbb{N}}$ converges to $x$.
    \end{itemize}
  \end{proof}
\end{proposition}

\begin{definition}
  Let $X$ be a set and $\set{x_n}_{n\in\mathbb{N}}\subseteq X$ be a sequence in $X$. A subsequence of $\set{x_n}_{n\in\mathbb{N}}$ is a sequence $\set{y_n}_{n\in\mathbb{N}}$ for which there exists a strictly increasing function $\psi:\mathbb{N}\rightarrow\mathbb{N}$ such that $y_n=x_{\varphi\pr{n}}$ for all $n\in\mathbb{N}$. We will denote $\set{y_n}_{n\in\mathbb{N}}:=\set{x_{n_k}}_{k\in\mathbb{N}}$ where $n_k:=\psi\pr{k}$ for all $k\in\mathbb{N}$. 
\end{definition}

\begin{remark}
  Let $X$ be a set and $\set{x_n}_{n\in\mathbb{N}}\subseteq X$ be a sequence in $X$. If $\set{y_{n_k}}_{k\in\mathbb{N}}$ is a subsequence of $\set{x_n}_{n\in\mathbb{N}}$, then $n_k<n_{k+1}$ for all $k\in\mathbb{N}$.
\end{remark}

\begin{proposition}
  Let $X$ be a set, $\set{x_n}_{n\in\mathbb{N}}\subseteq X$ be a sequence in $X$ and $x\in X$. Then $\set{x_n}_{n\in\mathbb{N}}$ converges to $x$ if and only if, for every subsequence $\set{x_{n_k}}_{k\in\mathbb{N}}$ of $\set{x_n}_{n\in\mathbb{N}}$ is convergent to $x$.
  \begin{proof}[\rm\textbf{Proof}]
    \hfill
    \begin{itemize}
      \item[$\boxed{\Rightarrow}$] Let $\set{x_{n_k}}_{k\in\mathbb{N}}$ be a subsequence of $\set{x_n}_{n\in\mathbb{N}}$ and $N\in\mathcal{N}_x\pr{X,\mathcal{T}}$ be a neighborhood of $x$. Then there exists $n_0:=n_0\pr{N,x}\in\mathbb{N}$ such that $x_n\in N$ for all $n\in\mathbb{N}$ with $n\geq n_0$. Take $k_0=k_0\pr{N,x}\in\mathbb{N}$ such that $n_{k_0}\geq n_0$. Then $n_k>n_0$ for all $k\in\mathbb{N}$ with $k\geq k_0$ and, consequently, $x_{n_k}\in N$ for all $k\in\mathbb{N}$ with $k\geq k_0$. Therefore, $\set{x_{n_k}}_{k\in\mathbb{N}}$ converges to $x$.
      \item[$\boxed{\Leftarrow}$] Trivial, because $\set{x_n}_{n\in\mathbb{N}}$ is a subsequence of itself.
    \end{itemize}
  \end{proof} 
\end{proposition}

\begin{proposition}
  Let $\pr{X,\mathcal{T}}$ be a topological space, $A\subseteq X$, $\set{x_n}_{n\in\mathbb{N}}\subseteq X$ be a sequence in $X$ convergent to $x\in X$. The following properties hold:

  \begin{enumerate}
    \item Let $\set{y_n}_{n\in\mathbb{N}}\subseteq X$ a sequence in $X$. Suppose there exists $n_1\in\mathbb{N}$ such that $y_n=x_n$ for all $n\in\mathbb{N}$ with $n\geq n_1$. Then $\set{y_n}_{n\in\mathbb{N}}$ converges to $x$.
    \item If $\set{x_n}_{n\in\mathbb{N}}\subseteq A$, then $x\in\overline{A}$. Particulary, if $A$ is closed, then $x\in A$.
    \item If $\set{x_n}_{n\in\mathbb{N}}\subseteq A\setminus\set{x}$, then $x\in A'$.
    \item Suppose that for every $x'\in X$, there exists a sequence $\set{x'_n}_{n\in\mathbb{N}}\subseteq A$ convergent to $x'$. Then $A$ is dense with respect to $\pr{X,\mathcal{T}}$.
    \item If $\set{x_n}_{n\in\mathbb{N}}\subseteq A$ is a sequence in $A$ and $\set{y_n}_{n\in\mathbb{N}}\subseteq X\setminus A$ is a sequence in $X\setminus A$ converging to $x$ as well, then $x\in\bd\pr{A}$.
  \end{enumerate}

  \begin{proof}[\rm\textbf{Proof}]
    \hfill
    \begin{enumerate}
      \item Let $N\in\mathcal{N}_x\pr{X,\mathcal{T}}$ be a neighborhood of $x$. Then there exists $n_2\in\mathbb{N}$ such that $x_n\in N$ for all $n\in\mathbb{N}$ with $n\geq n_2$. Take $n_0:=\max\set{n_1,n_2}\in\mathbb{N}$. Then $y_n=x_n\in N$ for all $n\geq n_0$. Therefore, $\set{y_n}_{n\in\mathbb{N}}$ converges to $x$.
      \item Let $N\in\mathcal{N}_x\pr{X,\mathcal{T}}$ be a neighborhood of $x$. Then there exists $n_0\in\mathbb{N}$ such that $x_n\in N$ for all $n\in\mathbb{N}$ with $n\geq n_0$. Since $\set{x_n}_{n\in\mathbb{N}}$ is a sequence in $A$, then $x_n\in A\cap N$ for all $n\in\mathbb{N}$ with $n\geq n_0$ and, consequently, $A\cap N\neq\varnothing$. Therefore, $x\in\overline{A}$.
      \item Since $\set{x_n}_{n\in\mathbb{N}}$ is sequence in $A\setminus\set{x}$ converging to $x$, by the previous property, we have $x\in\overline{A\setminus\set{x}}$. Therefore, $x\in A'$.
      \item By the first property, for all $x'\in X$, $x'\in\overline{A}$. Then $\overline{A}=X$. Therefore, $A$ is dense with respect to $\pr{X,\mathcal{T}}$.
      \item By the first property, since $\set{x_n}_{n\in\mathbb{N}}$ is a sequence in $A$ and $\set{y_n}_{n\in\mathbb{N}}$ is a sequence in $X\setminus A$ both converging to $x$, we have that $x\in\overline{A}$ and $x\in\overline{X\setminus A}$. Thus $x\in\overline{A}\cap\overline{X\setminus A}=\bd\pr{A}$.
    \end{enumerate}
  \end{proof}
\end{proposition}

\section{Nets}

\begin{definition}
  
\end{definition}
\chapter{Continuity}

\section{Continuous maps}

\begin{definition}
  Let $\pr{X,\mathcal{T}_1}$ and $\pr{Y,\mathcal{T}_2}$ be two topological spaces and $B\subseteq Y$. A map $f:X\rightarrow B$ is said to be \textit{$\pr{\mathcal{T}_1,\mathcal{T}_2}$-continuous} if $f^{-1}\pr{G}\in\mathcal{T}_1$ for any open set $G\in\mathcal{T}_2$. 
\end{definition}

\begin{definition}
  Let $\pr{X,\mathcal{T}_1}$ and $\pr{Y,\mathcal{T}_2}$ be two topological spaces and $B\subseteq Y$. A map $f:X\rightarrow B$ is said to be \textit{$\pr{\mathcal{T}_1,\mathcal{T}_2}$-continuous at $x\in X$} if for any neighborhood $N\in\mathcal{N}_{f\pr{x}}\pr{Y,\mathcal{T}_2}$ of $f\pr{x}$, there exists a neighborhood $M\in\mathcal{N}_x\pr{X,\mathcal{T}}$ of $x$ such that $f\pr{M}\subseteq N$.
\end{definition}

\begin{proposition}
  Let $\pr{X,\mathcal{T}_1}$ and $\pr{Y,\mathcal{T}_2}$ be two topological spaces with bases $\mathcal{B}_1\subseteq\mathcal{T}_1$ for $\mathcal{T}_1$ and $\mathcal{B}_2\subseteq\mathcal{T}_2$ for $\mathcal{T}_2$, $x\in X$, $B\subseteq Y$ and $f:X\rightarrow B$ be a map. The following are equivalent:
  
  \begin{enumerate}
    \item $f$ is  $\pr{\mathcal{T}_1,\mathcal{T}_2}$-continuous at $x$.
    \item For any open set $G\in\mathcal{T}_2$ containing $f\pr{x}$, there exists an open set $H=H\pr{x,G}\in\mathcal{T}_1$ containing $x$ such that $f\pr{G}\subseteq H$.
    \item For any basis set $B\in\mathcal{T}_2$ containing $f\pr{x}$, there exists a basis set $\widetilde{B}=\widetilde{B}\pr{x,G}\in\mathcal{T}_1$ containing $x$ such that $f\pr{\widetilde{B}}\subseteq B$.
  \end{enumerate}
  \begin{proof}[\rm\textbf{Proof}]
    
  \end{proof}
\end{proposition}

\begin{proposition}
  Let $\pr{X,\mathcal{T}_1}$ and $\pr{Y,\mathcal{T}_2}$ be two topological spaces, $a\in X'$, $B\subseteq Y$ and $f:X\rightarrow B$ be a map. Then $f$ is $\pr{\mathcal{T}_1,\mathcal{T}_2}$-continuous at $a$ if and only if $\lim_{x\to a}{f\pr{x}}=f\pr{a}$.
\end{proposition}

\begin{proposition}
  Let $\pr{X,\mathcal{T}_1}$ and $\pr{Y,\mathcal{T}_2}$ be two topological spaces, $B\subseteq Y$ and $f:X\rightarrow B$ be a map. The following are equivalent:

  \begin{enumerate}
    \item $f$ is $\pr{\mathcal{T}_1,\mathcal{T}_2}$-continuous.
    \item For every $x\in X$, $f$ is $\pr{\mathcal{T}_1,\mathcal{T}_2}$-continuous at $x$.
    \item If $F\in\mathcal{F}\pr{Y,\mathcal{T}_2}$ is closed, then $f^{-1}\pr{F}\in\mathcal{F}\pr{X,\mathcal{T}_1}$ is closed.
    \item If $A\subseteq X$, then $f\pr{\overline{A}}\subseteq\overline{f\pr{A}}$.
  \end{enumerate}
\end{proposition}

\subsection{Weak topology. Product topology}

\section{Open and closed maps. Homeomorphisms}

\begin{definition}
  Let $\pr{X,\mathcal{T}_1}$ and $\pr{Y,\mathcal{T}_2}$ be two topological spaces and $f:X\rightarrow Y$ be a map. We say that:

  \begin{itemize}
    \item $f$ is an \textit{open map} if $f\pr{G}\in\mathcal{T}_2$ is open for all $G\in\mathcal{T}_1$.
    \item $f$ is a \textit{closed map} if $f\pr{F}\in\mathcal{F}\pr{Y,\mathcal{T}_2}$ is closed for all $F\in\mathcal{F}\pr{X,\mathcal{T}_1}$.
  \end{itemize}
\end{definition}

\begin{definition}
  Let $\pr{X,\mathcal{T}_1}$ and $\pr{Y,\mathcal{T}_2}$ be two topological spaces and $f:X\rightarrow Y$ be a map. We say that $f$ is a \textit{homeomorphism} if $f$ is bijective and $\pr{\mathcal{T}_1,\mathcal{T}_2}$-continuous, and if its inverse is $\pr{\mathcal{T}_2,\mathcal{T}_1}$-continuous.
\end{definition}

\begin{proposition}
  Let $\pr{X,\mathcal{T}_1}$ and $\pr{Y,\mathcal{T}_2}$ be two topological spaces and $f:X\rightarrow Y$ be a bijective map. The following are equivalent:

  \begin{enumerate}
    \item $f$ is a homeomorphism.
    \item $f$ is open.
    \item $f$ is closed.
  \end{enumerate}

  \begin{proof}[\rm\textbf{Proof}]
    
  \end{proof}
\end{proposition}
\chapter{Separation and countability axioms}

\section{Separation axioms}

\begin{definition}
  A topological space $\pr{X,\mathcal{T}}$ is \textit{$T_0$ or Kolmogorov} if for any pair of points $x,y\in X$, there exists a neighborhood $N\in\mathcal{N}_x\pr{X,\mathcal{T}}$ of $x$ such that $y\notin N$ or there exists a neighborhood $M\in\mathcal{N}_y\pr{X,\mathcal{T}}$ of $y$ such that $x\notin M$.
\end{definition}

\begin{proposition}
  Let $\pr{X,\mathcal{T}}$ be a topological space and $\mathcal{B}\subseteq\mathcal{T}$ a basis for $\mathcal{T}$. The following are equivalent:

  \begin{enumerate}
    \item $\pr{X,\mathcal{T}}$ is Kolmogorov.
    \item For every pair of distinct points $x,y\in X$, there exists an open set $G\in\mathcal{T}$ such that $x\in G$ and $y\notin G$, or, $x\notin G$ and $y\in G$.
    \item For every pair of distinct points $x,y\in X$, there exists a basis set $B\in\mathcal{B}$ such that $x\in B$ and $y\notin B$, or, $x\notin B$ and $y\in B$.
  \end{enumerate}

  \begin{proof}[\rm\textbf{Proof}]
    \hfill
    \begin{itemize}
      \item[\boxed{1\Rightarrow2}]
      \item[\boxed{2\Rightarrow3}]
      \item[\boxed{3\Rightarrow1}]
    \end{itemize}
  \end{proof}
\end{proposition}

\begin{proposition}
  The following properties hold:

  \begin{enumerate}
    \item If $\pr{X,\mathcal{T}}$ is a Kolmogorov space and $A\subseteq X$, the subspace $\pr{A,\mathcal{T}_A}$ is a Kolmogorov space.
    \item If $\set{\pr{X_i,\mathcal{T}_i}}$ is an indexed collection of Kolmogorov spaces, then $\pr{\prod_{i\in I}{X_i},\mathcal{T}_{\prod_{i\in I}{X_i}}}$ is a Kolmogorov space.
  \end{enumerate}
  
  \begin{proof}[\rm\textbf{Proof}]
    \hfill
    \begin{enumerate}
      \item 
      \item 
    \end{enumerate}
  \end{proof}
\end{proposition}

\begin{definition}
  A topological space $\pr{X,\mathcal{T}}$ is \textit{$T_1$ or Fréchet} if for any pair of distinct points $x,y\in X$, there exists a neighborhood $N_1\in\mathcal{N}_x\pr{X,\mathcal{T}}$ of $x$ and a neighborhood $N_2\in\mathcal{N}_y\pr{X,\mathcal{T}}$ of $y$ such that $x\in N_1\setminus N_2$ and $y\in N_2\setminus N_1$.
\end{definition}

\begin{proposition}
  Let $\pr{X,\mathcal{T}}$ be a topological space and $\mathcal{B}\subseteq\mathcal{T}$ a basis for $\mathcal{T}$. The following are equivalent:

  \begin{enumerate}
    \item $\pr{X,\mathcal{T}}$ is Fréchet.
    \item For every pair of distinct points $x,y\in X$, there exists a pair of open sets $G_1,G_2\in\mathcal{T}$ such that $x\in G_1\setminus G_2$ and $y\in G_2\setminus G_1$.
    \item For every pair of distinct points $x,y\in X$,  there exists a pair of basis sets $B_1,B_2\in\mathcal{T}$ such that $x\in B_1\setminus B_2$ and $y\in B_2\setminus B_1$.
  \end{enumerate}
   
  \begin{proof}[\rm\textbf{Proof}]
    \hfill
    \begin{itemize}
      \item[\boxed{1\Rightarrow2}]
      \item[\boxed{2\Rightarrow3}]
      \item[\boxed{3\Rightarrow1}]
    \end{itemize}
  \end{proof}
\end{proposition}

\begin{proposition}
  A Fréchet space $\pr{X,\mathcal{T}}$ is a Kolmogorov space.
  \begin{proof}[\rm\textbf{Proof}]
    
  \end{proof}
\end{proposition}

\begin{proposition}
  The following properties hold:

  \begin{enumerate}
    \item If $\pr{X,\mathcal{T}}$ is a Fréchet space and $A\subseteq X$, the subspace $\pr{A,\mathcal{T}_A}$ is a Fréchet space.
    \item If $\set{\pr{X_i,\mathcal{T}_i}}$ is an indexed collection of Fréchet spaces, then $\pr{\prod_{i\in I}{X_i},\mathcal{T}_{\prod_{i\in I}{X_i}}}$ is a Fréchet space.
  \end{enumerate}
  
  \begin{proof}[\rm\textbf{Proof}]
    \hfill
    \begin{enumerate}
      \item 
      \item 
    \end{enumerate}
  \end{proof}
\end{proposition}

\begin{proposition}
  Let $\pr{X,\mathcal{T}}$ be a topological space. The following are equivalent:

  \begin{enumerate}
    \item $\pr{X,\mathcal{T}}$ is a Fréchet space.
    \item For every $x\in X$, $\set{x}\in\mathcal{F}\pr{X,\mathcal{T}}$ is a closed set.
    \item For every finite subset $A\subseteq X$ of $X$, $A\in\mathcal{F}\pr{X,\mathcal{T}}$ is a closed set.
  \end{enumerate}

  \begin{proof}[\rm\textbf{Proof}]
    \hfill
    \begin{enumerate}
      \item 
      \item 
      \item 
    \end{enumerate}
  \end{proof}
\end{proposition}

\begin{proposition}
  Let $\pr{X,\mathcal{T}}$ be a Fréchet space and $A\subseteq X$. Then $x\in A'$ if and only if for every neighborhood $N\in\mathcal{N}_x\pr{X,\mathcal{T}}$, the set $N\cap A$ has infinite points.
  
  \begin{proof}[\rm\textbf{Proof}]
    
  \end{proof}
\end{proposition}

\begin{definition}
  A topological space $\pr{X,\mathcal{T}}$ is \textit{$T_2$ or Hausdorff} if for any pair of distinct points $x,y\in X$, there exists a neighborhood $N_1\in\mathcal{N}_x\pr{X,\mathcal{T}}$ of $x$ and a neighborhood $N_2\in\mathcal{N}_y\pr{X,\mathcal{T}}$ such that $N_1\cap N_2=\varnothing$.
\end{definition}

\begin{proposition}
  Let $\pr{X,\mathcal{T}}$ be a topological space and $\mathcal{B}\subseteq\mathcal{T}$ a basis for $\mathcal{T}$. The following are equivalent:

  \begin{enumerate}
    \item $\pr{X,\mathcal{T}}$ is Hausdorff.
    \item For every pair of distinct points $x,y\in X$, there exists a pair of disjoint open sets $G_1,G_2\in\mathcal{T}$ such that $x\in G_1$ and $y\in G_2$.
    \item For every pair of distinct points $x,y\in X$, there exists a pair of disjoint basis sets $B_1,B_2\in\mathcal{T}$ such that $x\in B_1$ and $y\in B_2$.
  \end{enumerate}
  
  \begin{proof}[\rm\textbf{Proof}]
    \hfill
    \begin{itemize}
      \item[\boxed{1\Rightarrow2}]
      \item[\boxed{2\Rightarrow3}]
      \item[\boxed{3\Rightarrow1}]
    \end{itemize}
  \end{proof}
\end{proposition}

\begin{proposition}
  A Hausdorff space $\pr{X,\mathcal{T}}$ is a Fréchet space.

  \begin{proof}[\rm\textbf{Proof}]
    
  \end{proof}
\end{proposition}

\begin{proposition}
  If $\pr{X,\mathcal{T}}$ is a Hausdorff space and $\set{x_n}_{n\in\mathbb{N}}\subseteq X$ is a convergent sequence in $X$, there exists a unique $x\in X$ such that $\lim_{n\to\infty}{x_n}=x$.

  \begin{proof}[\rm\textbf{Proof}]
    
  \end{proof}
\end{proposition}

\begin{proposition}
  A topological space $\pr{X,\mathcal{T}}$ is Hausdorff if and only if for every convergent net $\set{x_\lambda}_{\lambda}\subseteq X$ in $X$, there exists a unique $x\in X$ such that $\lim_{\lambda\in\Lambda}{x_{\lambda}}=x$.

  \begin{proof}[\rm\textbf{Proof}]
    \hfill
    \begin{enumerate}
      \item[\boxed{\Rightarrow}]
      \item[\boxed{\Leftarrow}]
    \end{enumerate}
  \end{proof}
\end{proposition}

\begin{proposition}
  A Fréchet-Urysohn space $\pr{X,\mathcal{T}}$ is Hausdorff if and only if for every convergent sequence $\set{x_n}_{n\in\mathbb{N}}\subseteq X$ in $X$, there exists a unique $x\in X$ such that $\lim_{n\to\infty}{x_n}=x$.

  \begin{proof}[\rm\textbf{Proof}]
    \hfill
    \begin{enumerate}
      \item[\boxed{\Rightarrow}]
      \item[\boxed{\Leftarrow}]
    \end{enumerate}
  \end{proof}
\end{proposition}

\begin{proposition}
  A topological space $\pr{X,\mathcal{T}}$ is Hausdorff if and only if the diagonal $\Delta_X:=\set{\pr{x,x}\middle|x\in X}\in\mathcal{F}\pr{X\times X,\mathcal{T}_{X\times X}}$ is closed.
  
  \begin{proof}[\rm\textbf{Proof}]
    \hfill
    \begin{enumerate}
      \item[\boxed{\Rightarrow}]
      \item[\boxed{\Leftarrow}]
    \end{enumerate}
  \end{proof}
\end{proposition}

\begin{proposition}
  Let $\pr{X,\mathcal{T}_1}$ and $\pr{Y,\mathcal{T}_2}$ be topological spaces, $D\subseteq X$ be a dense set and $f,g:X\rightarrow Y$ two $\pr{\mathcal{T}_1,\mathcal{T}_2}$-continuous maps such that $f\pr{x}=g\pr{x}$ for all $x\in D$. Then $f=g$.

  \begin{proof}[\rm\textbf{Proof}]
    
  \end{proof}
\end{proposition}

\begin{proposition}
  Let $\pr{X,\mathcal{T}_1}$ and $\pr{Y,\mathcal{T}_2}$ be two topological spaces and $f:X\rightarrow Y$ be a map. If $\pr{Y,\mathcal{T}_2}$ is Hausdorff and $f$ is $\pr{\mathcal{T}_1,\mathcal{T}_2}$-continuous, then $\pr{X,\mathcal{T}}$ is Hausdorff.

  \begin{proof}[\rm\textbf{Proof}]
    \hfill
    \begin{enumerate}
      \item[\boxed{\Rightarrow}]
      \item[\boxed{\Leftarrow}]
    \end{enumerate}
  \end{proof}
\end{proposition}

\begin{proposition}
  The following properties hold:

  \begin{enumerate}
    \item If $\pr{X,\mathcal{T}}$ is a Hausdorff space and $A\subseteq X$, the subspace $\pr{A,\mathcal{T}_A}$ is a Hausdorff space.
    \item If $\set{\pr{X_i,\mathcal{T}_i}}$ is an indexed collection of Hausdorff spaces, then $\pr{\prod_{i\in I}{X_i},\mathcal{T}_{\prod_{i\in I}{X_i}}}$ is a Hausdorff space.
  \end{enumerate}
  
  \begin{proof}[\rm\textbf{Proof}]
    \hfill
    \begin{enumerate}
      \item 
      \item 
    \end{enumerate}
  \end{proof}
\end{proposition}

\begin{definition}
  A topological space $\pr{X,\mathcal{T}}$ is \textit{regular} if for any $x\in X$ and any closed set $F\in\mathcal{F}\pr{X,\mathcal{T}}$ not containing $x$, there exists a pair of disjoint open sets $G,H\in\mathcal{T}$ such that $x\in G$ and $F\subseteq H$.

  A $T_1$ and regular space is called a \textit{$T_3$ space}.
\end{definition}

\begin{proposition}
  A $T_3$ space $\pr{X,\mathcal{T}}$ is a Hausdorff space.
  \begin{proof}[\rm\textbf{Proof}]
    
  \end{proof}
\end{proposition}

\section{Countability axioms}


\input{General Topology/Topological spaces/Connectedness/connectedness}
\input{General Topology/Topological spaces/Compactness/compactness}
\input{General Topology/Metric spaces/metric_spaces.tex}

  \part{Measure Theory}
  \chapter{Introduction to Measure Theory}

\section{Measurable spaces}

\begin{definition}
  Let $X$ be a set. A \emph{$\sigma$-algebra on $X$} is a collection $\Sigma\subseteq\mathcal{P}\pr{X}$ of subsets of $X$ satisfying the following conditions:

  \begin{itemize}
    \item $\varnothing\in\Sigma$.
    \item If $A\in\Sigma$, then   $A^c:=X\setminus A\in\Sigma$.
    \item If $\set{A_n}_{n\in\mathbb{N}}\subseteq\Sigma$ is a countable collection of sets in $\Sigma$, then $\bigcup_{n=1}^{\infty}{A_n}\in\Sigma$.
  \end{itemize}

  The ordered pair $\pr{X,\Sigma}$ is called \emph{measurable space}.
\end{definition}

\section{Measure spaces}

\begin{definition}
  \rm Let $\pr{X,\Sigma}$ be a measurable space. A \emph{(positive) measure on $\pr{X,\Sigma}$} is a map $\mu:\Sigma\rightarrow\br{0,+\infty}$ satisfying the following conditions:

  \begin{itemize}
    \item $\mu\pr{\varnothing}=0$.
    \item If $\set{A_n}_{n\in\mathbb{N}}\subseteq\Sigma$ is a countable collection of pairwise disjoint sets in $\Sigma$, then $\mu\pr{\bigcup_{n=1}^{\infty}{A_n}}=\sum_{n=1}^{\infty}{\mu\pr{A_n}}$.
  \end{itemize}

  The triple $\pr{X,\Sigma,\mu}$ is called \emph{measure space}.
\end{definition}

\section{Measurable functions}

\begin{definition}
  \rm Let $\pr{X,\Sigma_1}$ and $\pr{Y,\Sigma_2}$ be two measurable spaces. We say a map $f:X\rightarrow Y$ is \emph{$\pr{\Sigma_1,\Sigma_2}$-measurable} if $f^{-1}\pr{B}\in\Sigma_1$ for all $B\in\Sigma_2$.
\end{definition}

\begin{proposition}
  \rm Let $\pr{X,\Sigma_1}$ and $\pr{Y,\Sigma_2}$ be two measurable spaces, $\mathcal{B}\subseteq\mathcal{P}\pr{Y}$ a set such that $\Sigma_2=\sigma\pr{Y,\mathcal{B}}$ and $f:X\rightarrow Y$ a map. Then $f$ is $\pr{\Sigma_1,\Sigma_2}$ measurable if and only if $f^{-1}\pr{B}\in\Sigma_1$ for all $B\in\mathcal{B}$.
\end{proposition}

\begin{definition}
  \rm Let $\pr{X,\Sigma}$ be a measurable space and let $\mathbb{K}\in\set{\mathbb{R},\mathbb{C}}$. A map $f:X\rightarrow\mathbb{K}$ is said to be \emph{$\Sigma$-measurable} if $f^{-1}\pr{G}\in\Sigma$ for every open set $G\subseteq\mathbb{K}$ (with respect to the Euclidean topology), or equivalently, $f$ is $\pr{\Sigma,\mathcal{B}\pr{\mathbb{K}}}$-measurable.
\end{definition}

\begin{proposition}
  Let $\pr{X,\Sigma}$ be a measurable space and $f:X\rightarrow\mathbb{C}$ a map. Then $f$ is $\Sigma$-measurable if and only if $\Re\pr{f},\Im\pr{f}:X\rightarrow\mathbb{R}$ are $\Sigma$-measurable.
\end{proposition}

\section{Lebesgue integral}

\section{Convergence theorems}

  \part{Functional Analysis}
  \input{Functional Analysis/Normed spaces/normed_spaces.tex}
\input{Functional Analysis/Pre-Hilbert spaces/pre_hilbert_spaces.tex}
\end{document}